% ClimateLens — คู่มือทดลองใช้และติดตั้ง (Group 1)
% Build:  xelatex ClimateLens-Guide.tex   (run twice for figure numbers)
% If private.tex exists (gitignored), the grader login and our Supabase keys are filled in;
% otherwise the PDF says where to find them. make_private.py writes private.tex.
\documentclass[11pt,a4paper]{article}

\usepackage{fontspec}
\usepackage[a4paper,margin=2.3cm,footskip=1.2cm]{geometry}
\usepackage{float,graphicx,xcolor,enumitem,fancyhdr,caption,booktabs,array,tabularx}
\usepackage{fvextra}
\usepackage[most]{tcolorbox}
\usepackage[hidelinks]{hyperref}

\setmainfont{Sarabun}[
  Path=fonts/, Extension=.ttf,
  UprightFont=*-Regular, BoldFont=*-Bold,
  ItalicFont=*-Italic, BoldItalicFont=*-BoldItalic,
  Scale=1.05]
\setmonofont{Latin Modern Mono}[Scale=1.0]
\XeTeXlinebreaklocale "th"
\XeTeXlinebreakskip = 0pt plus 1pt
\linespread{1.25}
\setlength{\parindent}{0pt}
\setlength{\parskip}{6pt}

\definecolor{accent}{HTML}{E8763A}
\definecolor{codebg}{HTML}{F5F5F2}
\definecolor{notebg}{HTML}{FFF4EC}

\pagestyle{fancy}
\fancyhf{}
\renewcommand{\headrulewidth}{0pt}
\fancyfoot[C]{\small ClimateLens · Group 1 · 06016474 Applied AI for IT \quad \thepage}

\captionsetup{font={it,small},labelsep=space}
\renewcommand{\figurename}{รูปที่}
\renewcommand{\tablename}{ตารางที่}
\renewcommand{\contentsname}{สารบัญ}

\setlist[enumerate,1]{label=\arabic*.,leftmargin=*,itemsep=3pt}
\setlist[enumerate,2]{label=\alph*.,itemsep=1pt}
\setlist[itemize]{itemsep=2pt}

% Section titles in the style of the course labs: "ส่วนที่ 1: ..."
\makeatletter
\renewcommand\section{\@startsection{section}{1}{0pt}{14pt plus 4pt}{6pt}{\Large\bfseries}}
\renewcommand\subsection{\@startsection{subsection}{2}{0pt}{10pt plus 3pt}{4pt}{\large\bfseries}}
\renewcommand\thesection{ส่วนที่ \arabic{section}:}
\renewcommand\thesubsection{\arabic{section}.\arabic{subsection}}
\makeatother

% Code blocks: ASCII only. Macros (\GraderEmail ...) expand inside them.
\DefineVerbatimEnvironment{code}{Verbatim}{breaklines,fontsize=\small,commandchars=\\\{\}}
\tcbset{codebox/.style={colback=codebg,colframe=codebg,boxrule=0pt,arc=2pt,left=6pt,right=6pt,top=3pt,bottom=3pt}}
\newenvironment{codeblock}{\begin{tcolorbox}[codebox]}{\end{tcolorbox}}
\newtcolorbox{note}{colback=notebg,colframe=accent,boxrule=0pt,leftrule=3pt,arc=0pt,left=8pt,top=4pt,bottom=4pt}

\newcommand{\ui}[1]{\textbf{#1}}          % a button or menu name
\newcommand{\file}[1]{\texttt{#1}}
\newcommand{\fig}[3][0.85]{%              % \fig[width]{file}{caption}
  \begin{figure}[H]\centering
  \IfFileExists{figures/#2}{\fbox{\includegraphics[width=#1\linewidth]{figures/#2}}}%
    {\fbox{\parbox[c][3cm][c]{#1\linewidth}{\centering (ภาพ #2)}}}%
  \caption{#3}\end{figure}}

% ----- Values that must not be in the public repo -----
\newif\ifprivate
\IfFileExists{private.tex}{\input{private.tex}\privatetrue}{%
  \newcommand{\GraderEmail}{\textrm{(ดูในข้อความที่ส่งใน Google Classroom)}}
  \newcommand{\GraderPassword}{\textrm{(ดูในข้อความที่ส่งใน Google Classroom)}}
  \newcommand{\SupaURL}{https://<project-ref>.supabase.co}
  \newcommand{\SupaKey}{sb\_publishable\_...}}

\newcommand{\AppURL}{http://ito.it.kmitl.ac.th:8501}
\newcommand{\AppIP}{http://10.70.38.15:8501}
\newcommand{\Typhoon}{hf.co/typhoon-ai/typhoon2.5-qwen3-4b-gguf:Q4\_K\_M}

\begin{document}

\begin{center}
  {\LARGE\bfseries ClimateLens}\\[2pt]
  {\large คู่มือทดลองใช้งานและติดตั้งระบบ}\\[6pt]
  {\small Group 1 — Anawat Wongprachanukul 66070217 · Amorntep Thongplum 66070219 ·
   Pubeth Sriwattana 66070158 · Teerapat Tuanpromrat 66070093}
\end{center}

ClimateLens รับความคิดเห็นปลายเปิดภาษาไทยจากแบบสำรวจพนักงาน (CSV) แล้ว
(1) จำแนก\ui{ความรู้สึก} (Positive / Neutral / Negative)
(2) จัด\ui{หัวข้อ} 7 หมวดด้วย LLM ที่รันบนเครื่องเอง
(3) คำนวณ \ui{Climate Index} แยกตามหัวข้อ แผนก และเดือน แล้วแสดงบน Dashboard

เลือกวิธีทดลองได้ 3 แบบ เรียงจากเร็วที่สุด:

\begin{center}
\begin{tabularx}{\linewidth}{@{}l X l@{}}
\toprule
\textbf{วิธี} & \textbf{ต้องมี} & \textbf{เวลา} \\
\midrule
ส่วนที่ 2: ใช้เว็บที่ deploy แล้ว & เบราว์เซอร์ + เครือข่ายคณะ IT & 2 นาที \\
ส่วนที่ 3: รันบนเครื่องตัวเอง (ใช้ฐานข้อมูลของกลุ่ม) & Python 3.12, Ollama & 20–30 นาที \\
ส่วนที่ 4: ติดตั้งทั้งระบบ (Supabase ของผู้ทดสอบเอง) & + บัญชี Supabase & 45–60 นาที \\
\bottomrule
\end{tabularx}
\end{center}

% =====================================================================
\section{ภาพรวมระบบและผลของโมเดล}

\begin{enumerate}
  \item \ui{Upload CSV} — คอลัมน์ \file{feedback} (บังคับ), \file{department}, \file{date} (\file{YYYY-MM-DD})
  \item \ui{Sentiment} — Character TF-IDF (2–5 ตัวอักษร) + Logistic Regression
        เทรนบน Wisesight Sentiment Corpus (\file{models/sentiment.joblib})
  \item \ui{Topic} — Typhoon 2.5 (Qwen3 4B) แบบ zero-shot รันผ่าน Ollama
        ความคิดเห็นไม่ถูกส่งไป LLM บนคลาวด์
  \item \ui{Save} — เก็บลง Supabase (PostgreSQL) แยกข้อมูลรายผู้ใช้ด้วย Row Level Security
  \item \ui{Dashboard} — Climate Index $= 50 + 50\,(P-N)/T$;
        $\ge 60$ Healthy, 45–60 Watch, $<45$ Needs attention; กลุ่มที่มีน้อยกว่า 10 ความเห็นไม่แสดงสถานะ
\end{enumerate}

\begin{table}[h!]\centering\small
\caption{ผลการประเมินโมเดล (ตัวเลขเดียวกับสไลด์นำเสนอ)}
\begin{tabular}{@{}llll@{}}
\toprule
\textbf{งาน} & \textbf{ชุดทดสอบ} & \textbf{Accuracy} & \textbf{Macro F1} \\
\midrule
Sentiment (โมเดลที่ใช้) & Wisesight test 2,614 แถว & 69.7\% & 0.668 \\
Sentiment baseline (ตอบ Neutral ทุกแถว) & Wisesight test 2,614 แถว & 55.6\% & 0.238 \\
Topic — Typhoon 4B zero-shot (ที่ใช้ในแอป) & ข้อสอบ 105 ข้อ (7 หมวด × 15) & 86.7\% & 0.867 \\
Topic — Typhoon 4B few-shot & ข้อสอบ 105 ข้อ & 87.6\% & 0.886 \\
\bottomrule
\end{tabular}
\end{table}

\begin{note}
\textbf{ข้อจำกัดที่ทราบ:} โมเดล Sentiment เทรนจากโพสต์โซเชียลมีเดีย แต่ใช้กับข้อความในที่ทำงาน
(Positive F1 = 0.52) จึงมักอ่านความเห็นเชิงบวกเป็น Neutral — จะเห็นได้ในไฟล์ทดลองด้านล่าง
\end{note}

% =====================================================================
\section{ทดลองบนเว็บที่ deploy แล้ว}

\begin{note}
เว็บอยู่บน VM ของคณะ เข้าได้จาก\textbf{เครือข่ายคณะ IT / VPN ของ ITKMITL เท่านั้น}
\end{note}

\begin{enumerate}
  \item เปิด \url{http://ito.it.kmitl.ac.th:8501} (หรือ \url{http://10.70.38.15:8501})
  \item เข้าสู่ระบบด้วยบัญชีทดลอง
    \begin{enumerate}
      \item Email: \texttt{\GraderEmail}
      \item Password: \texttt{\GraderPassword}
    \end{enumerate}
  \item บัญชีนี้มีข้อมูลตัวอย่าง 258 ความเห็น (พ.ค.–ต.ค. 2026, 3 แผนก) ดูได้ทันทีที่หน้า
        \ui{Overview}, \ui{Departments}, \ui{Trends}, \ui{Feedback Explorer}
  \item ทดลองอัปโหลด: ไปที่ \ui{Upload \& Analyze} แล้วเลือกไฟล์จากโฟลเดอร์ \file{demo/try/}
    \begin{enumerate}
      \item \file{climatelens\_grader\_5.csv} — 5 แถว ประมาณ 20 วินาที
      \item \file{climatelens\_demo\_live\_10.csv} — 10 แถว ประมาณ 40 วินาที
    \end{enumerate}
  \item ตรวจผลในตาราง แล้วกด \ui{Save to Supabase} (ไม่บังคับ) จากนั้นกด
        \ui{Use this data in the dashboard} เพื่อดูไฟล์ที่เพิ่งอัปโหลดบน Dashboard
\end{enumerate}

\fig[0.5]{app01-login.png}{หน้าเข้าสู่ระบบ}
\fig[0.72]{app02-overview.png}{หน้า Overview: Climate Index แยกแผนก × หัวข้อ}
\fig[0.72]{app03-upload.png}{หน้า Upload \& Analyze หลังวิเคราะห์ไฟล์ 5 แถว}

\begin{note}
\textbf{เรื่องที่ควรสังเกต:} Climate Index รวมของ 3 เดือนล่าสุดคือ 53.9 (ดูปกติ) แต่ช่อง Production × Workload ได้ 31 (สีแดง)
— นี่คือเหตุผลที่ต้องแยก Climate Index ตามหัวข้อ
\end{note}

% =====================================================================
\section{รันบนเครื่องตัวเอง (ใช้ฐานข้อมูลของกลุ่ม)}

ใช้ Supabase ของกลุ่ม จึงไม่ต้องตั้งค่าฐานข้อมูล เข้าสู่ระบบด้วยบัญชีทดลองเดียวกับส่วนที่ 2

\subsection{ติดตั้งโปรแกรม}
\begin{enumerate}
  \item ติดตั้ง Python 3.12 (\url{https://www.python.org/downloads/})
  \item ติดตั้ง Ollama (\url{https://ollama.com/download}) แล้วเปิดโปรแกรมทิ้งไว้
  \item ดาวน์โหลดโมเดล Typhoon (ประมาณ 2.5 GB, ใช้ RAM ราว 4 GB):
\begin{codeblock}\begin{code}
ollama pull hf.co/typhoon-ai/typhoon2.5-qwen3-4b-gguf:Q4_K_M
\end{code}
\end{codeblock}
  \item แตกไฟล์ \file{group\_01.zip} หรือ clone repository:
\begin{codeblock}\begin{code}
git clone https://github.com/AAnawat/Applied-AI-Project.git
cd Applied-AI-Project
\end{code}
\end{codeblock}
  \item สร้าง virtual environment และติดตั้ง package (รันที่โฟลเดอร์หลักของโปรเจกต์)
\begin{codeblock}\begin{code}
python -m venv .venv
.venv\textbackslash{}Scripts\textbackslash{}activate        (Windows)
source .venv/bin/activate      (macOS / Linux)
pip install -r requirements.txt
\end{code}
\end{codeblock}
\end{enumerate}

\subsection{ตั้งค่าการเชื่อมต่อ}
\begin{enumerate}
  \item คัดลอก \file{.streamlit/secrets.toml.example} เป็น \file{.streamlit/secrets.toml}
  \item แก้สองบรรทัดแรกเป็นค่าของกลุ่ม:
\begin{codeblock}\begin{code}
supabase_url = "\SupaURL"
supabase_key = "\SupaKey"
\end{code}
\end{codeblock}
\ifprivate\else ค่าจริงของทั้งสองบรรทัดอยู่ในข้อความที่ส่งใน Google Classroom\fi
\end{enumerate}

\begin{note}
\file{supabase\_key} นี้คือ \textbf{Publishable key} ซึ่ง Supabase ออกแบบให้เปิดเผยได้ (ฝังอยู่ในหน้าเว็บทั่วไป)
ข้อมูลถูกป้องกันด้วย Row Level Security — ผู้ใช้แต่ละคนเห็นเฉพาะข้อมูลของตัวเอง
ส่วน \textbf{Secret key} ไม่ได้แจกให้ และแอปจะปฏิเสธถ้าใส่ Secret key
\end{note}

\subsection{เปิดแอป}
\begin{enumerate}
  \item รันจาก\textbf{โฟลเดอร์หลักของโปรเจกต์} (ไม่ใช่ใน \file{dashboard/}):
\begin{codeblock}\begin{code}
python -m streamlit run dashboard/app.py
\end{code}
\end{codeblock}
  \item เปิด \url{http://localhost:8501} แล้วเข้าสู่ระบบด้วยบัญชีทดลอง
  \item ทดลองอัปโหลดตามข้อ 4–5 ของส่วนที่ 2
\end{enumerate}

% =====================================================================
\section{ติดตั้งทั้งระบบด้วย Supabase ของตัวเอง}

ใช้เมื่อต้องการฐานข้อมูลแยกของตัวเองทั้งหมด ทำส่วนที่ 3.1 ให้เสร็จก่อน

\subsection{สร้างโปรเจกต์ Supabase}
\begin{enumerate}
  \item สร้างโปรเจกต์ใหม่ตาม Lab 12 Section 1 (\ui{New project}, Region: Singapore, เปิด \ui{Enable Data API})
\end{enumerate}
\fig{sb01-new-project.png}{ปุ่ม New project ในหน้า Projects}

\subsection{สร้างตารางด้วย SQL}
\begin{enumerate}
  \item ไปที่เมนู \ui{SQL Editor} แล้วกด \ui{+ New}
  \item เปิดไฟล์ต่อไปนี้ทีละไฟล์ \textbf{ตามลำดับ} คัดลอกทั้งไฟล์ไปวาง แล้วกด \ui{Run}
    \begin{enumerate}
      \item \file{supabase/schema.sql} — ตาราง, view และสิทธิ์
      \item \file{supabase/migrations/20261025000000\_auth\_rls.sql} — เจ้าของข้อมูล + Row Level Security
      \item \file{supabase/migrations/20261025120000\_survey\_date.sql} — วันที่ของแบบสำรวจ
    \end{enumerate}
  \item แต่ละไฟล์ต้องขึ้น \texttt{Success. No rows returned}
\end{enumerate}
\fig{sb02-sql-editor.png}{SQL Editor: วาง SQL แล้วกด Run}

\begin{note}
ไฟล์ที่ 2 มีตัวแปร \file{owner\_email} ไว้ย้ายข้อมูลเก่า — โปรเจกต์ใหม่ยังไม่มีข้อมูล \textbf{ไม่ต้องแก้}
\end{note}

\subsection{ตั้งค่าการเข้าสู่ระบบ}
\begin{enumerate}
  \item ไปที่ \ui{Authentication} → \ui{Sign In / Providers}
  \item เปิด \ui{Allow new users to sign up} ถ้าต้องการให้สมัครเองได้จากแท็บ \ui{Register} ในแอป
  \item ปิด \ui{Confirm email} (ระบบส่งอีเมลของ Supabase ส่งได้ไม่กี่ฉบับต่อชั่วโมง) แล้วกด \ui{Save changes}
\end{enumerate}
\fig{sb05-auth-signups.png}{Authentication → Sign In / Providers}

\subsection{สร้างผู้ใช้}
\begin{enumerate}
  \item ไปที่ \ui{Authentication} → \ui{Users} → \ui{Add user} → \ui{Create new user}
  \item กรอก Email และ Password, ติ๊ก \ui{Auto confirm user?} แล้วกด \ui{Create user}
\end{enumerate}
\fig{sb06-add-user-menu.png}{เมนู Add user → Create new user}
\fig[0.4]{sb07-create-user.png}{ฟอร์มสร้างผู้ใช้ (ติ๊ก Auto confirm user)}

\subsection{คัดลอก URL และ Publishable key}
\begin{enumerate}
  \item ไปที่ \ui{Project Settings} → \ui{Data API} คัดลอก \ui{API URL}
        แล้ว\textbf{ตัด \file{/rest/v1/} ท้ายออก} ให้เหลือ \file{https://<project-ref>.supabase.co}
  \item ไปที่ \ui{Project Settings} → \ui{API Keys} คัดลอก \ui{Publishable key}
        (\file{sb\_publishable\_...}) — \textbf{ห้ามใช้ Secret key}
  \item ใส่ทั้งสองค่าใน \file{.streamlit/secrets.toml} ตามส่วนที่ 3.2 แล้วเปิดแอปตามส่วนที่ 3.3
\end{enumerate}
\fig{sb04-data-api-url.png}{Data API: API URL}
\fig{sb03-api-keys.png}{API Keys: ใช้ Publishable key เท่านั้น}

\subsection{ใส่ข้อมูลตัวอย่าง}
\begin{enumerate}
  \item เข้าสู่ระบบในแอปด้วยผู้ใช้ที่สร้างในข้อ 4.4
  \item หน้า \ui{Upload \& Analyze}: อัปโหลดไฟล์ใน \file{demo/history/} ทีละเดือน (พ.ค. → ต.ค.)
        แล้วกด \ui{Save to Supabase} ทุกไฟล์
  \item ใช้\textbf{ไฟล์รายเดือน 6 ไฟล์ หรือ ไฟล์รวม 258 แถว อย่างใดอย่างหนึ่ง} — ถ้าใช้ทั้งคู่ข้อมูลจะซ้ำ
\end{enumerate}

\begin{note}
เวลาวิเคราะห์หัวข้อด้วย CPU: ประมาณ 1.3–3.6 วินาที/แถว (ไฟล์รายเดือน 43 แถว ≈ 1–2.5 นาที)
ถ้ารีเฟรชหน้าระหว่างวิเคราะห์ จะต้องเริ่มไฟล์นั้นใหม่
\end{note}

% =====================================================================
\section{Deploy บนเซิร์ฟเวอร์ด้วย Docker}

ใช้ \file{docker-compose.yml} รันแอป + Ollama บนเครื่องเดียว (ทดสอบบน Ubuntu 24.04, 8 vCPU, 8 GB RAM, ไม่มี GPU)

\begin{enumerate}
  \item คัดลอกโปรเจกต์ไปที่เซิร์ฟเวอร์ และสร้าง \file{.streamlit/secrets.toml} ตามส่วนที่ 3.2
  \item เพิ่มบรรทัดนี้ใน \file{secrets.toml} ให้แอปคุยกับ container ของ Ollama:
\begin{codeblock}\begin{code}
ollama_host = "http://ollama:11434"
\end{code}
\end{codeblock}
  \item รันที่โฟลเดอร์หลักของโปรเจกต์ (ครั้งแรกจะดาวน์โหลดโมเดล 2.5 GB):
\begin{codeblock}\begin{code}
docker compose up -d --build
docker compose ps
\end{code}
\end{codeblock}
  \item \file{docker compose ps} ต้องแสดง \file{app} เป็น \file{healthy}
  \item เปิด \texttt{http://<server-ip>:8501}
\end{enumerate}

% =====================================================================
\section{ปัญหาที่พบบ่อย}

\begin{center}\small
\begin{tabularx}{\linewidth}{@{}>{\raggedright\arraybackslash}p{0.44\linewidth} >{\raggedright\arraybackslash}X@{}}
\toprule
\textbf{ข้อความที่เห็น} & \textbf{วิธีแก้} \\
\midrule
\texttt{Set supabase\_url and supabase\_key in .streamlit/secrets.toml}
  & ยังไม่ได้สร้าง \file{secrets.toml} หรือรันจากโฟลเดอร์ \file{dashboard/} — รันจากโฟลเดอร์หลัก \\\addlinespace[3pt]
\texttt{The configured Supabase key is a Secret/service\_role key}
  & ใส่ Secret key ผิดช่อง — เปลี่ยนเป็น Publishable key \\\addlinespace[3pt]
\texttt{Wrong email or password.}
  & ตรวจบัญชี หรือยังไม่ได้ติ๊ก Auto confirm ตอนสร้างผู้ใช้ \\\addlinespace[3pt]
\texttt{Could not analyze topics with Local LLM}
  & Ollama ไม่ได้เปิด หรือยังไม่ได้ \texttt{ollama pull} โมเดล \\\addlinespace[3pt]
\texttt{Missing `feedback` column}
  & CSV ต้องมีหัวคอลัมน์ \file{feedback} และบันทึกเป็น UTF-8 \\\addlinespace[3pt]
\texttt{Please remove result columns before analyzing}
  & ลบคอลัมน์ \file{sentiment} / \file{category} ออกจาก CSV \\\addlinespace[3pt]
บันทึกไม่ได้ (code \texttt{42501})
  & ยังไม่ได้รัน SQL ไฟล์ที่ 2 (\file{auth\_rls}) \\\addlinespace[3pt]
บันทึกไม่ได้ (code \texttt{PGRST204})
  & ยังไม่ได้รัน SQL ไฟล์ที่ 3 (\file{survey\_date}) \\
\bottomrule
\end{tabularx}
\end{center}

\end{document}
